% ch07_u07 —— 原书 p.296--p.305（7.10.20 证明续 -- 7.11.6）
% 本单元自 7.10.20 证明的续文起（承上单元 p.295 末句），至 7.11.6 末句止；
% 尾部接 p.306。

\SourcePage{296}

逆向态射把 $(\mathcal{R},\beta)$ 送到由平凡 gerbe 借助
$(\mathcal{R},\beta)$ 粘成的对 $(\mathcal{C},\mathcal{E}_V)$。具体地说，
$(\mathcal{C},\mathcal{E}_V)$ 的定义如下。由于 $V$ 形式光滑，结构态射
$j:S\to X$ 可延拓为 $j':\hat S\to V$。此时 $\mathcal{C}(\hat S)$ 是一个
$\mathcal{P}ic(\hat S)$-torsor，并带有如下附加结构：

\begin{enumerate}[label=(\roman*),leftmargin=2.2em]
\item 对任一如上之 $j'$，给定 $\mathcal{C}(\hat S)$ 的一个对象，记作
$j'^*\mathcal{E}_V$。
\item 若 $j'':\hat S\to V$ 是 $j$ 的另一延拓，则我们有一个线丛之间的等同
$\theta_{j''j'}:\operatorname{Hom}(j'^*\mathcal{E}_V,j''^*\mathcal{E}_V)
\xrightarrow{\ \sim\ }(j'',j')^*\mathcal{R}$。
\end{enumerate}

我们要求 (ii) 把 $\operatorname{Hom}$ 的复合与 $\beta$ 所定义的同构等同起来。
容易看出，这样的 $\mathcal{C}(\hat S)$ 存在且唯一（相差唯一的等价）。诸纤维
$\mathcal{C}(\hat S)$ 以显然的方式粘成一个晶体 $\mathcal{O}^*$-gerbe。按构造有
$\mathcal{E}_V\in\mathcal{C}(V)$。\hfill$\square$

\bdp{7.10.21} \emph{注.} 设 $\mathcal{E}'_V$ 是 $\mathcal{C}(V)$ 的另一对象，
而 $(\mathcal{R}',\beta')\in G(V)$ 是对应于 $(\mathcal{C},\mathcal{E}'_V)$ 的对。
置 $\mathcal{L}:=\operatorname{Hom}(\mathcal{E}_V,\mathcal{E}'_V)
\in\mathcal{P}ic(V)$。则
\[
\mathcal{R}'=\operatorname{Ad}_{\mathcal{L}}\mathcal{R}
:=(p_2^*\mathcal{L})\otimes\mathcal{R}\otimes(p_1^*\mathcal{L})^{\otimes-1},
\qquad
\beta'=\operatorname{Ad}_{\mathcal{L}}\beta .
\]
现在设 $\mathcal{C}$ 是 $X$ 上任一晶体 $\mathcal{O}^*$-gerbe，并设我们有如上的
$X\hookrightarrow V$。要用 7.10.20 描述 $\mathcal{C}$，必须保证
$\mathcal{C}(V)$ 非空。

% [Tutor Note: 原书 p.296 此处正文只写「$\mathcal{C}(V)$ is non-empty
% if$^{*)}\ldots$」，脚注 \*) 仅作「and, certainly, only if」；正文的 if 与脚注
% 合读才是「当且仅当」。译文按原书拆分照录，不合并。]
\bdp{7.10.22} \emph{引理.} 设 $X$ 仿射，且 $V$ 是可数多个子概形之并。则
$\mathcal{C}(V)$ 非空，只要 $\mathcal{C}(X,X)$ 非空%
\footnote{并且当然，仅当。}。

\emph{证明.} 设 $X'\subset V$ 是 $X$ 的一个无穷小邻域。则任一
$\mathcal{E}_X\in\mathcal{C}(X,X)$ 都可延拓为
$\mathcal{E}_{X'}\in\mathcal{C}(X,X')$，且所有这样的延拓彼此同构。现在有 $X$
的无穷小邻域序列 $X\subset X^{(1)}\subset X^{(2)}\cdots$，使得
$V=\varinjlim X^{(n)}$。归纳地定义序列
$\mathcal{E}_X(n)\in\mathcal{C}(X,X^{(n)})$ 以及等同
$\mathcal{E}_X(n+1)|_{X^{(n)}}=\mathcal{E}_X(n)$。这就是
$\mathcal{E}_V\in\mathcal{C}(V)$。\hfill$\square$

\SourcePage{297}

\bdp{7.10.23} \emph{注.} (i) 考虑 $X_{\text{ét}}$ 上的 $\mathcal{O}^*$-gerbe
$i\mathcal{C}$（故 $i\mathcal{C}(U)=\mathcal{C}(U,U)$）。则
$\mathcal{C}(X,X)\neq\varnothing$ 当且仅当 $i\mathcal{C}$ 是平凡 gerbe，即它在
$H^2(X_{\text{ét}},\mathcal{O}^*)=Br(X)$ 中的类消没%
\footnote{这个类是 $\mathcal{C}$ 的类在映射
$H^2(X_{cr},\mathcal{O}^*_{cr})\to
H^2(X_{cr},i\cdot\mathcal{O}^*)=Br(X)$ 下的像。}。

(ii) 对任意代数空间 $X$ 与 $\mathcal{C}\in\mathcal{G}_{cr}(X)$，可在
$X_{\text{ét}}$ 局部用 7.10.20 描述 $\mathcal{C}$。具体地说，存在 $X$ 的平展
覆盖 $U_i$，使 $U_i$ 仿射且 $\mathcal{C}(U_i,U_i)\neq\varnothing$。把 $U_i$
嵌入光滑概形，并取相应的形式完备化 $V_i$。于是由 7.10.22，可用 7.10.20、
7.10.21 描述 $\mathcal{C}_{U_i}$。

\bdp{7.10.24} \emph{定义.} 设 $\mathcal{C}\in\mathcal{G}_{cr}(X)$。$X$ 上的
$\mathcal{C}$-扭 $\mathcal{D}$-晶体是一个 Cartesian 函子
$M:\mathcal{C}\to\mathcal{M}^{!}(X_{cr},\mathcal{O})$，使得对任一
$\varepsilon\in\mathcal{C}(\hat S)$ 与 $f\in\mathcal{O}^*(\hat S)$，有
$M(f\varepsilon)=f\cdot\operatorname{id}_{M(\varepsilon)}$。

$\mathcal{C}$-扭 $\mathcal{D}$-晶体构成一个 $\mathbb{C}$-范畴
$\mathcal{M}_{\mathcal{C}}(X)$。它以函子性的方式依赖于 $\mathcal{C}$
（对态射 $\mathcal{C}\to\mathcal{C}'$，相应地有范畴等价
$\mathcal{M}_{\mathcal{C}}(X)\xrightarrow{\ \sim\ }
\mathcal{M}_{\mathcal{C}'}(X)$，等等）。

诸范畴 $\mathcal{M}_{\mathcal{C}}(U)=\mathcal{M}_{\mathcal{C}_U}(U)$，
$U\in X_{\text{ét}}$，以显然的方式构成 $X_{\text{ét}}$ 上的范畴层
$\mathcal{M}_{\mathcal{C}}(X_{\text{ét}})$。

设 $\mathcal{C}^{triv}$ 是平凡化 gerbe，故
$\mathcal{C}^{triv}(\hat S)=\mathcal{P}ic(\hat S)$。以
$\mathcal{C}^{triv}$ 扭的 $\mathcal{D}$-晶体与通常的 $\mathcal{D}$-晶体相同。
具体地说，把 $M\in\mathcal{M}_{\mathcal{C}^{triv}}(X)$ 与 $\mathcal{D}$-晶体
$M_{\hat S}:=M(\mathcal{O}_{\hat S})$ 等同起来。

\emph{注.} 在上述定义中，可把 $X_{cr}$ 换成
$X_{\text{ét}\,cr}$。若 $X$ 是概形且
$\mathcal{C}\in\mathcal{G}_{Zar\,cr}(X)$，则可把 $X_{cr}$ 换成 $X_{Zar\,cr}$。
得到的是同一个范畴 $\mathcal{M}_{\mathcal{C}}(X)$。

\bdp{7.10.25} 这里是 7.10.7、7.10.8 的扭版本。设我们处于情形 7.10.19，即
有 $(\mathcal{C},\mathcal{E}_V)\in\mathcal{G}^V_{cr}(X)$ 以及相应的
$(\mathcal{R},\beta)\in G(V)$（见 7.10.20）。范畴
$\mathcal{M}_{\mathcal{C}}(X)$ 可描述如下。设
$\mathcal{M}_{\mathcal{R}}(X)=\mathcal{M}_{\mathcal{R}\beta}(X)$ 是二元对
$(M_V,\tau)$ 的范畴，其中 $M_V\in\mathcal{M}(V,\mathcal{O})$，而
$\tau:(p_1^{!}M_V)\otimes\mathcal{R}\xrightarrow{\ \sim\ }p_2^{!}M_V$
是 $\mathcal{M}(V^{<2>},\mathcal{O})$\SourcePage{298}中的同构，且满足%
\footnote{我们用 $\beta$ 来等同等式左右两端所在的模。}
\begin{equation}\tag{344}
p_{23}^{!}(\tau)p_{12}^{!}(\tau)=p_{13}^{!}(\tau).
\end{equation}

\bdp{7.10.26} \emph{引理.} 范畴 $\mathcal{M}_{\mathcal{C}}(X)$ 与
$\mathcal{M}_{\mathcal{R}}(X)$ 典范等价。

\emph{证明.} 对 $M\in\mathcal{M}_{\mathcal{C}}(X)$，置
$M_V=M(\mathcal{E}_V):=\bigcup M(\mathcal{E}_{(X,X')})$，并把 $\tau$ 定义为诸同构
\[
(p_1^{!}M_V)\otimes\mathcal{R}=M(p_1^*\mathcal{E}_V)\otimes\mathcal{R}
=M((p_1^*\mathcal{E}_V)\otimes\mathcal{R})
=M(p_2^*\mathcal{E}_V)=p_2^{!}M_V
\]
的复合。其余部分是 7.10.8 的证明的直接改写。\hfill$\square$

\bdp{7.10.27} \emph{引理.} 对任意 $X$ 与 $\mathcal{C}\in\mathcal{G}_{cr}(X)$，
范畴 $\mathcal{M}_{\mathcal{C}}(X)$ 是阿贝尔范畴。

\emph{证明.} 对 7.10.9 的证明作显然的修改。用到 7.10.23(ii)、7.10.22、
7.10.26。\hfill$\square$

\bdp{7.10.28} 从现在起，我们设 $X$ 是光滑代数空间。我们要把上述图景与 tdo
以及扭 $\mathcal{D}$-模的通常框架作比较。首先，把晶体
$\mathcal{O}^*$-gerbe 与 tdo 联系起来\footnote{关于 tdo 的基本事实，例如参看
[BB93] 2.1。}。

考察 $V=X$ 时的 7.10.19。考虑由对 $(\mathcal{C},\mathcal{E}_X)$ 构成的
Picard 群胚 $\mathcal{G}^{\sim}_{cr}(X):=\mathcal{G}^{V}_{cr}(X)$，其中
$\mathcal{C}$ 是 $X$ 上的晶体 $\mathcal{O}^*$-gerbe，而
$\mathcal{E}_X\in\mathcal{C}(X)$。

$\mathcal{G}^{\sim}_{cr}(X)$ 有一个方便的解释。考虑 $X$ 上的
$\mathcal{I}_{cr}$-gerbe（即 $X_{cr}$ 上的 $\mathcal{I}_{cr}$-gerbe）。由于
$H^0(X_{cr},\mathcal{I}_{cr})=0$，这些 gerbe 构成（平移的）Picard 群胚
$\mathcal{G}\mathcal{I}_{cr}(X)$。指数态射
$\mathcal{I}_{cr}\hookrightarrow\mathcal{O}^*_{cr}$ 给出函子
$\exp:\mathcal{G}\mathcal{I}_{cr}(X)\to\mathcal{G}_{cr}(X)$。由于
$\mathcal{I}_{(X,X)}=0$，对任一 $\mathcal{I}_{cr}$-gerbe $\mathcal{B}$，群胚
$\mathcal{B}_X$ 平凡，故群胚 $(\exp\mathcal{B})_X$ 有特异对象
$\mathcal{E}_{\mathcal{B}X}$（在典范同构意义下确定）。这样我们就定义了 Picard
群胚的态射
\begin{equation}\tag{345}
\exp:\mathcal{G}\mathcal{I}_{cr}(X)\longrightarrow \mathcal{G}^{\sim}_{cr}(X),
\end{equation}

\SourcePage{299}

$\mathcal{B}\mapsto(\exp\mathcal{B},\mathcal{E}_{\mathcal{B}X})$。这是 Picard
群胚的等价（这由 (342) 推出）。

\emph{例.} (342) 的「边界映射」给出 Picard 群胚的态射
$c:\mathcal{P}ic(X)\to\mathcal{G}\mathcal{I}_{cr}(X)$（即晶体 Chern 类）。
用 (345) 的术语，它把 $\mathcal{L}\in\mathcal{P}ic(X)$ 送到对
$(\mathcal{C}^{triv},\mathcal{L})$。

\bdp{7.10.29} \emph{命题.} $\mathcal{G}^{\sim}_{cr}(X)$ 与 $X$ 上 tdo 构成的
Picard 群胚 $\mathcal{T}\mathcal{D}\mathcal{O}(X)$ 典范等价。

\emph{证明.} 按 $V=X$ 时的 7.10.20，把我们的 $\mathcal{G}^{\sim}_{cr}(X)$ 与
$G(X)$ 等同起来。现设 $(\mathcal{R},\beta)\in G(X)$，相应的 tdo
$\mathcal{D}_{\mathcal{R}}=\mathcal{D}_{(\mathcal{R},\beta)}$ 如下定义。我们沿用
7.10.13 的记号。把 $\mathcal{D}_X$ 看作 Diff-双模（即
$\mathcal{M}(X^{<2>},\mathcal{O})$ 的一个对象）。置
\[
\mathcal{D}_{\mathcal{R}}:=
\mathcal{D}_X\otimes_{\mathcal{O}_{X^{<2>}}}\mathcal{R}.
\]
乘法态射
$m_{\mathcal{R}}:\mathcal{D}_{\mathcal{R}}\otimes_{\mathcal{O}_X}
\mathcal{D}_{\mathcal{R}}\to p_{13}^{!}\mathcal{D}_{\mathcal{R}}$ 是
$\mathcal{D}_X$ 的相应态射与 $\beta$ 的张量积。容易验证
$\mathcal{D}_{\mathcal{R}}$ 是 tdo，且 $G(X)\to\mathcal{T}\mathcal{D}\mathcal{O}(X)$、
$(\mathcal{R},\beta)\mapsto\mathcal{D}_{\mathcal{R}}$ 是 Picard 群胚的态射。

逆向态射把 $X$ 上的 tdo $A$ 送到对象 $(\mathcal{R},\beta)$，其中
$\mathcal{R}:=\operatorname{Hom}_{\mathcal{O}_{X^{<2>}}}(\mathcal{D}_X,A)$，而
$\beta$ 由比较 $\mathcal{D}_X$ 与 $A$ 的乘法态射 $m$ 来定义。细节留给读者。
\hfill$\square$

\bdp{7.10.30} \emph{注.} 上述等价还有另一种（等价的）写法。由 (345)，
$\mathcal{G}^{\sim}_{cr}(X)$ 等价于 $\mathcal{G}\mathcal{I}_{crys}(X)$，即等价于
与复形 $\tau_{\le1}(R\Gamma(X_{crys},\mathcal{I}_{Xcrys})[1])$ 相伴的 Picard
群胚。按 [BB93] 2.1.6、2.1.4，$\mathcal{T}\mathcal{D}\mathcal{O}(X)$ 是与复形
$\tau_{\le1}(R\Gamma(X,\Omega_X^{\ge1})[1])$ 相伴的 Picard 群胚。现在上述两个复形
典范拟同构（见 7.10.18(iii)）。

\bdp{7.10.31} 这里是 7.10.12 的扭版本。对
$(\mathcal{C},\mathcal{E}_X)\in\mathcal{G}^{\sim}_{cr}(X)$，考虑相应的
$(\mathcal{R},\beta)\in G(X)$ 与 tdo $\mathcal{D}_{\mathcal{R}}$。取
$M\in\mathcal{M}_{\mathcal{C}}(X)$。按 7.10.26，可把 $M$ 看作
$\mathcal{M}_{\mathcal{R}}(X)$ 中的对 $(M_X,\tau)$。\SourcePage{300}由于%
\footnote{参看 7.10.12 与 7.10.29 的证明。}
$p_1^{!}M_X=M_X\otimes_{\mathcal{O}_X}\mathcal{D}_X$ 且
$\mathcal{D}_{\mathcal{R}}=
\mathcal{D}_X\otimes_{\mathcal{O}_{X^{<2>}}}\mathcal{R}$，我们可以把 $\tau$
改写成 $\mathcal{M}(X^{<2>},\mathcal{O})$ 中的同构
\begin{equation}\tag{346}
M_X\otimes_{\mathcal{O}_X}\mathcal{D}_{\mathcal{R}}
\xrightarrow{\ \sim\ }p_2^{!}M_X .
\end{equation}
由伴随性，可把 (346) 看作 $\mathcal{O}_X$-模的态射
\begin{equation}\tag{347}
M_X\otimes\mathcal{D}_{\mathcal{R}}\to M_X .
\end{equation}
记 $\mathcal{M}^{r}(X,\mathcal{D}_{\mathcal{R}})$ 为 $X$ 上右
$\mathcal{D}_{\mathcal{R}}$-模的范畴。

\bdp{7.10.32} \emph{命题.} 态射 (347) 是 $\mathcal{D}_{\mathcal{R}}$ 在
$M_X$ 上的右幺作用。函子
$\mathcal{M}_{\mathcal{C}}(X)\to\mathcal{M}^{r}(X,\mathcal{D}_{\mathcal{R}})$、
$M\mapsto M_X$ 是范畴等价。

\emph{证明.} 留给读者（见 7.10.12、7.10.13）。\hfill$\square$

\bdsec{7.11. 归纳概形上的 $\mathcal{D}$-模}

本节讨论形式光滑归纳概形上的 $\mathcal{D}$-模理论。注意 $\mathcal{D}$-晶体
图景（见 7.10）在归纳概形的框架下立刻有意义；而通常的进路（微分算子等）
要费些篇幅才写得下来。

\bdp{7.11.1} 严格意义下的归纳概形 $X$ 是一个「空间」（即交换
$\mathbb{C}$-代数范畴上的集合值函子
$A\mapsto X(A)=X(\operatorname{Spec}A)$），它可以表为
$\varinjlim X_\alpha$，其中 $\{X_\alpha\}$ 是拟紧概形的有向族，使得所有映射
$i_{\alpha\beta}:X_\alpha\to X_\beta$（$\alpha\le\beta$）都是闭嵌入。若 $X$
能如上表出，且指标 $\alpha$ 的集合可数，则称 $X$ 是
$\aleph_0$-归纳概形\footnote{并非所有自然的归纳概形例子都是
$\aleph_0$-归纳概形；例如对每个无穷维向量空间 $V$，函子
$A\mapsto\operatorname{End}_A(V\otimes A)$ 是归纳概形，但不是
$\aleph_0$-归纳概形。}。若 $P$\SourcePage{301}是概形的一个在取闭子概形下
保持的性质，则称 $X$ 满足归纳-$P$ 性质，如果每个 $X_\alpha$ 都满足 $P$。

置 $X_{red}:=\varinjlim X_{\alpha red}$；若 $X_{red}=X$，则称归纳概形 $X$
是既约的。

形式概形是这样的归纳概形 $X$：$X_{red}$ 是概形（关于形式概形的这一定义与
EGA 中定义的关系，参看 7.12.17 中的讨论）。$\aleph_0$-形式概形是
$\aleph_0$-归纳概形的形式概形。归纳概形 $Z$ 沿闭子概形 $Y\subset Z$ 的
完备化，是满足 $Y'_{red}=Y_{red}$ 的闭子概形 $Y'\subset Z$ 的正向极限。
对形式概形，我们不写「归纳仿射」而写作「仿射」。形式概形 $X$ 仿射，当且仅当
$X_{red}$ 仿射。

仿照 Grothendieck（[Gr64]、[Gr67]），若对每个 $A$ 与每个幂零理想
$I\subset A$，映射 $X(A)\to X(A/I)$ 都是满射，则称 $X$ 形式光滑。容易看出，
对归纳有限型的归纳概形，形式光滑性是局部性质（参看 [Gr67] 中命题 17.1.6
的证明）\footnote{我们不知道对非归纳有限型的归纳概形这是否成立。对概形，
答案是肯定的。这可由 [Gr68a] 中的注 9.5.8 与下述出人意料的结果（[RG]，
p.~82，3.1.4）推出：投射模这一性质对 Zariski 拓扑是局部的，甚至对 fpqc
拓扑也是局部的（无需任何有限性假设！）。}。若对任意如上之 $A$、$I$，从
$X(A)$ 到 $Y(A)$ 与 $X(A/I)$ 在 $Y(A/I)$ 上的纤维积的映射是满射，则称
态射 $X\to Y$ 形式光滑。

设 $X$ 是归纳概形。若对任意满足 $Y\subset Z$ 的闭子概形 $Z\subset X$，$Y$
在 $\mathcal{O}_Z$ 中的理想都有限生成，则称闭拟紧子概形 $Y\subset X$ 是
合理的。若 $X$ 是其合理子概形之并，即 $X$ 可表为 $\varinjlim X_\alpha$
且所有 $X_\alpha$ 都合理，则称 $X$ 是合理的。合理归纳概形的闭子空间是合理
归纳概形；两个合理归纳概形之积合理。

\SourcePage{302}

\emph{注.} 把上述定义中的「概形」一词换成「代数空间」，就得到归纳代数空间
的概念。全部讨论都自动适用于归纳代数空间的情形。

\bdp{7.11.2} \emph{例.} (i) 归纳仿射的归纳概形 $X$ 等同于一个投射代数，即
交换代数范畴中的一个投射对象 $R$，它可表为 $\varprojlim R_\alpha$，使得映射
$R_\beta\to R_\alpha$（$\beta\ge\alpha$）都是满射。我们写
$X=\operatorname{Spf}R:=\varinjlim\operatorname{Spec}R_\alpha$。拓扑由开理想
$I_\alpha\subset R$ 定义的完备拓扑交换代数 $R$ 可看作投射代数（置
$R_\alpha:=R/I_\alpha$）。并非所有投射代数都是这种类型：若指标 $\alpha$ 的
集合不可数，则从 $R_\alpha$ 的集合论投射极限到 $R_{\alpha_0}$ 的映射未必是
满射\footnote{即使映射 $R_\beta\to R_\alpha$（$\beta\ge\alpha$）都是满射
（正如我们所设）。}。当然，归纳仿射的 $\aleph_0$-归纳概形，等同于拓扑由
$R$ 的可数或有限个开理想系定义的完备拓扑代数。

(ii) 设 $V$ 是 Tate 向量空间（见 4.2.13）。则 $V$（更确切地说，函子
$A\mapsto V\hat\otimes A$）是合理的归纳仿射归纳概形。事实上，$V$ 中每个
c-格都是仿射概形。有 $V=\operatorname{Spf}R$，其中
$R=\varprojlim\operatorname{Sym}(U_\alpha^*)$，$U_\alpha$ 跑遍 $V$ 中
c-格的集合。

若 $X$ 是合理归纳概形，则对 $x\in X(\mathbb{C})$，$X$ 在 $x$ 处的切空间
$\Theta_x$ 是 Tate 向量空间：$\Theta_x$ 的拓扑由 $X$ 的含 $x$ 的合理子概形
在 $x$ 处的切空间定义。于是若 $H$ 是合理群归纳概形，则其李代数
$\operatorname{Lie}H$ 是 Tate 向量空间范畴中的李代数。

(iii) 对如上之 $V$，记 $Gr(V)$ 为 $V$ 中 c-格的「空间」。更确切地说，
$Gr(V)$ 是把交换代数 $A$ 送到 $V\hat\otimes A$ 中 c-格（按 4.2.14 的意义）
的集合的函子。显然 $Gr(V)$ 是归纳固有、形式光滑的归纳概形（事实上，它是
所有 c-格对 $U_1\subset U_2\subset V$ 对应的 $U_2/U_1$ 的 Grassmann 概形
之并）。

\SourcePage{303}

(iv) 设 $K$ 是局部域，$O\subset K$ 是相应的局部环（故
$K\simeq\mathbb{C}((t))$，$O\simeq\mathbb{C}[[t]]$）。对任意「空间」$Y$，
我们有「空间」$Y(O)\subset Y(K)$，其定义为 $Y(O)(A):=
Y(A\hat\otimes O)$、$Y(K)(A):=Y(A\hat\otimes K)$（这里
$A\hat\otimes O=A[[t]]$，$A\hat\otimes K=A((t))$）。设 $Y$ 是仿射概形。则
$Y(O)$ 也是仿射概形，而 $Y(K)$ 是归纳仿射的 $\aleph_0$-归纳概形。若 $Y$
有限型，则 $Y(K)$ 合理。若 $Y$ 光滑，则 $Y(O)$ 与 $Y(K)$ 都形式光滑。

设 $G$ 是仿射代数群，$\mathfrak{g}$ 是其李代数。考虑群归纳概形 $G(K)$。
有 $\operatorname{Lie}(G(K))=\mathfrak{g}(K)=\mathfrak{g}\otimes K$，
$\operatorname{Lie}(G(O))=\mathfrak{g}(O)=\mathfrak{g}\otimes O$。

(v) 设 $G$ 是合理群归纳概形，且 $G_{red}$ 是仿射群概形。函子
$G\mapsto(\operatorname{Lie}G,G_{red})$ 是如上诸 $G$ 的范畴与 Harish-Chandra
对的范畴之间的等价。对归纳概形 $X$，$G$ 在 $X$ 上的作用，等同于
$(\operatorname{Lie}G,G_{red})$ 在 $X$ 上的作用。类似地，$G$-模等同于
$(\operatorname{Lie}G,G_{red})$-模，等等。

\bdp{7.11.3} 在归纳概形的框架下，有两种不同的方式定义 $\mathcal{O}$-模；
相应的对象分别称为 $\mathcal{O}^{p}$-模与 $\mathcal{O}^{!}$-模。我们从较直接
（但较不重要）的 $\mathcal{O}^{p}$-模概念开始\footnote{这里「$p$」代表
「投射极限」。}，它对任意「空间」$X$ 都有意义（见 7.11.1）。

$X$ 上的 $\mathcal{O}^{p}$-模 $P$ 是这样一条规则：它把交换代数 $A$ 与点
$\phi\in X(A)$ 送到 $A$-模 $P_\phi$，并把代数的态射 $f:A\to B$ 送到
$B$-模的等同 $f_P:B\otimes_f P_\phi\xrightarrow{\ \sim\ }P_{f\phi}$。这些等同
须与 $f$ 的复合相容。若 $X=\varinjlim X_\alpha$ 是归纳概形，则这样的 $P$
等同于：$X_\alpha$ 上（拟凝聚）$\mathcal{O}$-模 $P_{X_\alpha}$ 的族，连同对
$\alpha\le\beta$ 的等同 $i^*_{\alpha\beta}P_{X_\beta}=P_{X_\alpha}$，且满足
显然的传递性。若每个 $P_\phi$（或每个 $P_{X_\alpha}$）都平坦，则称 $P$ 平坦。
$X$ 上可逆 $\mathcal{O}^{p}$-模（又称线丛）的定义类似。

\SourcePage{304}

记 $X$ 上 $\mathcal{O}^{p}$-模的范畴为 $\mathcal{M}^{p}(X,\mathcal{O})$。这是
张量 $\mathbb{C}$-范畴。$\mathcal{M}^{p}(X,\mathcal{O})$ 中的单位对象是函数的
「层」$\mathcal{O}_X$。注意 $\mathcal{M}^{p}(X,\mathcal{O})$ 未必是阿贝尔
范畴。平坦 $\mathcal{O}^{p}$-模的范畴
$\mathcal{M}^{pfl}(X,\mathcal{O})$ 是（Quillen 意义下的）正合范畴。

对任意 $P,P'\in\mathcal{M}^{p}(X,\mathcal{O})$，向量空间
$\operatorname{Hom}(P,P')$ 带有显然的拓扑；态射的复合是连续的。特别地，
$\Gamma(X,P):=\operatorname{Hom}(\mathcal{O}_X,P)$ 是拓扑向量空间，且是拓扑环
$\Gamma(X,\mathcal{O}_X)$ 上的模。

\emph{注.} (i) 若把 $\mathcal{O}$-模换成任一在仿射概形范畴上纤维化的范畴，
上述定义仍有意义。例如可考虑左 $\mathcal{D}$-模（又称带可积联络的
$\mathcal{O}$-模）；归纳概形上相应的对象称为（左）$\mathcal{D}^{p}$-模。

(ii) 若 $X$ 是归纳仿射的 $\aleph_0$-归纳概形，
$X=\operatorname{Spf}R=\varinjlim\operatorname{Spec}R/I_\alpha$
（见 7.11.2(i)），则 $X$ 上的 $\mathcal{O}^{p}$-模等同于完备且分离的拓扑
$R$-模 $P$，使得 $I_\alpha P\subset P$ 的闭包构成拓扑的基。

\bdp{7.11.4} 下面转向 $\mathcal{O}^{!}$-模。这里必须设 $X$ 是合理归纳概形。
$X$ 上的 $\mathcal{O}^{!}$-模 $M$ 是这样一条规则：它把合理子概形 $Y\subset X$
送到拟凝聚 $\mathcal{O}_Y$-模 $M_{(Y)}$，并对 $Y\subset Y'$ 给出态射
$M_{(Y)}\to M_{(Y')}$。这些态射把 $M_{(Y)}$ 与
$i^{!}_{YY'}M_{(Y')}:=\operatorname{Hom}_{\mathcal{O}_{Y'}}
(\mathcal{O}_Y,M_{(Y')})$ 等同起来，且满足显然的传递性条件%
\footnote{我们需要考虑合理子概形，才能保证 $i^{!}$ 保持拟凝聚性。}。若写
$X=\varinjlim X_\alpha$ 且诸 $X_\alpha$ 合理，则只需考虑诸 $X_\alpha$，而不必
考虑所有合理子概形。$X$ 上的 $\mathcal{O}^{!}$-模构成阿贝尔范畴
$\mathcal{M}(X,\mathcal{O})$。注意对任意闭子概形 $Y\subset X$，范畴
$\mathcal{M}(Y,\mathcal{O})$ 是 $\mathcal{M}(X,\mathcal{O})$ 的全子范畴，且对
子商封闭；又对任一 $\mathcal{O}^{!}$-模 $M$，有
$M=\varinjlim M_{(X_\alpha)}$。

范畴 $\mathcal{M}(X,\mathcal{O})$ 是张量范畴
$\mathcal{M}^{p}(X,\mathcal{O})$ 上的模。具体地说，对
$M\in\mathcal{M}(X,\mathcal{O})$、$P\in\mathcal{M}^{p}(X,\mathcal{O})$，其张量
积 $M\otimes P\in$
\SourcePage{305}$\mathcal{M}(X,\mathcal{O})$ 是
$\varinjlim_\alpha M_{(X_\alpha)}
\otimes_{\mathcal{O}_{X_\alpha}}P_{X_\alpha}$。函子
$\otimes:\mathcal{M}(X,\mathcal{O})\times
\mathcal{M}^{pfl}(X,\mathcal{O})\to\mathcal{M}(X,\mathcal{O})$ 双正合。

对 $\mathcal{O}^{!}$-模 $M$，我们把其整体截面的空间定义为
$\Gamma(X,M):=\varinjlim_\alpha\Gamma(X_\alpha,M_{(X_\alpha)})$。函子
$\Gamma(X,\cdot)$ 左正合。

\emph{注.} (i) 诸范畴 $\mathcal{M}(Y,\mathcal{O})$ 连同函子
$i^{!}_{YY'}$ 一起，构成 $X$ 的合理子概形范畴（有序集）上的纤维化范畴；
而 $\mathcal{M}(X,\mathcal{O})$ 是其 Cartesian 截面的范畴。

(ii) 若 $X=\operatorname{Spf}R$ 且投射代数 $R$ 是拓扑代数（见 7.11.2），则
$X$ 上的 $\mathcal{O}^{!}$-模等同于离散 $R$-模（这里「离散」是指 $R$-作用
关于 $M$ 上的离散拓扑连续）。

(iii) 若 $P$ 平坦，则
$(M\otimes P)_{(X_\alpha)}=M_{(X_\alpha)}\otimes P_{X_\alpha}$。

\bdp{7.11.5} 设群归纳概形（或任意群「空间」）$K$ 作用在 $X$ 上。则对每个
交换代数 $A$，群 $K(A)$ 作用在 $\operatorname{Spec}A\times X$ 上。对
$M\in\mathcal{M}(X,\mathcal{O})$，$K$ 在 $M$ 上的作用由
$\mathcal{O}_{\operatorname{Spec}A}\boxtimes M\in
\mathcal{M}(\operatorname{Spec}A\times X,\mathcal{O})$ 上的 $K(A)$-作用给出，
且要求对任意态射 $A\to A'$ 相应的作用相容。记 $X$ 上 $K$-等变
$\mathcal{O}^{!}$-模的范畴为
$\mathcal{M}(K\backslash X,\mathcal{O})$。$K$-等变 $\mathcal{O}^{p}$-模的
定义留给读者。

\bdp{7.11.6} 7.10 的所有基本定义与结果（拓扑 $X_{cr}$、$\mathcal{D}$-晶体、
晶体 $\mathcal{O}^*$-torsor、扭 $\mathcal{D}$-晶体、基本函子性这些概念的定义）对任意
归纳有限型的归纳概形 $X$ 都显然有意义。于是从 $\mathcal{D}$-晶体的观点看，
$\mathcal{D}$-模理论自动推广到归纳概形的框架。

本节剩余部分要讨论的是 $\mathcal{D}$-模的通常进路（微分算子环等），它在我们
的归纳概形形式光滑时适用。比较 $\mathcal{D}$-晶体与 $\mathcal{D}$-模两种
框架的结果 7.10.12、7.10.29、7.10.32，对形式光滑的归纳概形逐字成立。

% 尾部接 p.306（本单元末句完结于 7.11.6 的「remain literally true for
% formally smooth ind-schemes」。）
